\subsection{Двойственность в пространствах Соболева. Пространства с отрицательными индексами}

\begin{remark}[с лекции]
Этот параграф дальше в курсе не используется; материал — 70-е годы XX века.
\end{remark}

\subsubsection*{Общий вид функционала на $\Wp{p}{\ell}$}

\begin{theorem}\label{thm:w-dual}
Пусть $1 < p < \infty$, $N = \sum_{|\alpha| \le \ell} 1$ и $f \in \big(\Wp{p}{\ell}(\Om)\big)^*$.
Тогда существует $\underline v = (v_\alpha)_{|\alpha| \le \ell} \in \big[\Lp{p^*}(\Om)\big]^N$
такой, что
\[
  f(u) = \sum_{|\alpha| \le \ell} \int\limits_\Om v_\alpha\, D^\alpha u\, dx
  \qquad \forall u \in \Wp{p}{\ell}(\Om) ,
\]
причём $\norm{\underline v}_{[\Lp{p^*}(\Om)]^N} = \norm{f}_{(\Wp{p}{\ell}(\Om))^*}$.
\end{theorem}

\begin{idea}
Это \emph{не} теорема Рисса, а более простой случай: вся работа уже сделана в предыдущем
параграфе, где $\Wp{p}{\ell}(\Om)$ отождествлено с подпространством
$\overline{W} \subset \big[\Lp{p}(\Om)\big]^N$. Остаётся продолжить функционал с
подпространства на всё $\big[\Lp{p}\big]^N$ и применить уже известное описание
сопряжённого к $\big[\Lp{p}\big]^N$.
\end{idea}

\begin{warn}
В тетради в формулировке дважды написано $\Wpz{p}{\ell}$ и $\Wp{p}{\ell}$ вперемешку.
Доказательство ведётся для полного $\Wp{p}{\ell}(\Om)$ (оператор $P$ определён на нём),
а переход к $\Wpz{p}{\ell}$ нужен лишь дальше, когда определяются пространства
с отрицательными индексами.
\end{warn}

\begin{proof}
\textbf{1. Оператор $P$.} Положим
\[
  P u = (D^\alpha u)_{|\alpha| \le \ell} \in \big[\Lp{p}(\Om)\big]^N ,
  \qquad
  \overline{W} = \Big\{ \underline w \in \big[\Lp{p}(\Om)\big]^N \ \Big|\
  \exists\, u \in \Wp{p}{\ell}(\Om): w_\alpha = D^\alpha u \Big\} .
\]
Как показано в предыдущем параграфе, $P$ — изометрический изоморфизм
$\Wp{p}{\ell}(\Om)$ на $\overline W$, а $\overline W$ — подпространство
в $\big[\Lp{p}(\Om)\big]^N$ (и потому рефлексивно и сепарабельно при $1 < p < \infty$).

\textbf{2. Перенос функционала на $\overline W$.} Положим
\[
  \widetilde f(\underline w) = f\big(P^{-1}\underline w\big), \qquad \underline w \in \overline W .
\]
Функционал $\widetilde f$ линеен (как композиция линейных $f$ и $P^{-1}$) и ограничен:
\[
  |\widetilde f(\underline w)| = \big| f(P^{-1}\underline w) \big|
  \le \norm{f}_{(\Wp{p}{\ell})^*} \norm{P^{-1}\underline w}_{\Wp{p}{\ell}(\Om)}
  = \norm{f}_{(\Wp{p}{\ell})^*} \norm{\underline w}_{[\Lp{p}(\Om)]^N} ,
\]
где использована изометричность $P$. Значит,
$\norm{\widetilde f}_{\overline W^{\,*}} \le \norm{f}$; обратно, подставляя
$\underline w = Pu$,
\[
  |f(u)| = |\widetilde f(Pu)| \le \norm{\widetilde f}\, \norm{Pu}_{[\Lp{p}]^N}
  = \norm{\widetilde f}\, \norm{u}_{\Wp{p}{\ell}(\Om)} ,
\]
то есть $\norm{f} \le \norm{\widetilde f}$ и $\norm{\widetilde f} = \norm{f}$.

\textbf{3. Продолжение и представление.} По теореме Хана--Банаха продолжим $\widetilde f$
с подпространства $\overline W$ на всё $\big[\Lp{p}(\Om)\big]^N$ без увеличения нормы:
получим $F \in \big(\big[\Lp{p}(\Om)\big]^N\big)^*$ с $\norm{F} = \norm{\widetilde f} = \norm{f}$.
По описанию сопряжённого к $\big[\Lp{p}\big]^N$ существует
$\underline v \in \big[\Lp{p^*}(\Om)\big]^N$ с
\[
  F(\underline w) = \sum_{|\alpha| \le \ell} \int\limits_\Om v_\alpha w_\alpha\, dx ,
  \qquad \norm{\underline v}_{[\Lp{p^*}]^N} = \norm{F} = \norm{f} .
\]
Подставляя $\underline w = Pu$, получаем требуемое представление.
\end{proof}

\begin{idea}
Вектор $\underline v$, в отличие от скалярного случая, \textbf{не единствен}: продолжение
по Хану--Банаху с подпространства $\overline W$ на всё $\big[\Lp{p}\big]^N$ неоднозначно,
и разные продолжения дают разные $\underline v$ с одним и тем же $f$. Единственной
остаётся только норма.
\end{idea}

\begin{idea}[title={Чему равно $N$}]
$N$ — число мультииндексов $\alpha \in \N_0^n$ с $|\alpha| \le \ell$. Добавив фиктивную
переменную $\alpha_{n+1} = \ell - |\alpha| \ge 0$, получаем биекцию с наборами
$(\alpha_1, \dots, \alpha_{n+1})$, в сумме дающими ровно $\ell$, поэтому
\[
  N = \binom{n + \ell}{n} .
\]
В тетради рядом с $N = \sum_{|\alpha|\le\ell} 1$ стоит вопрос «какая формула?» — вот она.
\end{idea}

\subsubsection*{Пространства $\Wp{p}{-\ell}(\Om)$}

Ограничим функционал на финитные функции. Для $\underline v \in \big[\Lp{p}(\Om)\big]^N$
распределение
\[
  T(\varphi) = \sum_{|\alpha| \le \ell} (-1)^{|\alpha|}\, D^\alpha T_{v_\alpha}(\varphi)
  = \sum_{|\alpha| \le \ell} \int\limits_\Om v_\alpha\, D^\alpha \varphi\, dx ,
  \qquad \varphi \in C_0^\infty(\Om) ,
\]
оценивается через $\norm{\varphi}_{\ell,p^*,\Om}$, то есть непрерывно по норме
$\Wp{p^*}{\ell}$. Так как $C_0^\infty(\Om)$ плотно в $\Wpz{p^*}{\ell}(\Om)$ по определению,
по теореме о продолжении ограниченного оператора с плотной области определения существует
единственный $f \in \big(\Wpz{p^*}{\ell}(\Om)\big)^*$ с
\[
  f(\varphi) = T(\varphi) \quad \forall \varphi \in C_0^\infty(\Om) ,
  \qquad
  \norm{f}_{(\Wpz{p^*}{\ell})^*} = \norm{T} .
\]

\begin{definition}
Распределение $T$ принадлежит $\Wp{p}{-\ell}(\Om)$, если существует
$\underline v \in \big[\Lp{p}(\Om)\big]^N$ такое, что
\[
  T(\varphi) = \sum_{|\alpha| \le \ell} (-1)^{|\alpha|} D^\alpha T_{v_\alpha}(\varphi)
  \qquad \forall \varphi \in C_0^\infty(\Om) .
\]
Норма:
\[
  \norm{T}_{\Wp{p}{-\ell}(\Om)}
  = \sup \Big\{ |T(\varphi)| \ \Big|\ \varphi \in C_0^\infty(\Om),\ \norm{\varphi}_{\ell,p^*,\Om} = 1 \Big\} .
\]
\end{definition}

\begin{warn}
В тетради показатель у $\Lp{\;}$ в определении читается неразборчиво. Выбран вариант,
согласованный с утверждением ниже: коэффициенты $v_\alpha$ берутся из $\Lp{p}$, и тогда
спаривание с $D^\alpha\varphi \in \Lp{p^*}$ корректно по Гёльдеру. При обратном выборе
показателей утверждение $\Wp{p}{-\ell} \sim \big(\Wpz{p^*}{\ell}\big)^*$ не сходится.
\end{warn}

\begin{proposition}
\[
  \Wp{p}{-\ell}(\Om) \ \sim\ \big(\Wpz{p^*}{\ell}(\Om)\big)^* ,
\]
то есть соответствие $T \leftrightarrow f$ выше — биекция, сохраняющая норму.
\end{proposition}

\begin{proof}
В одну сторону это только что построенное продолжение $T \mapsto f$ с сохранением нормы.
В другую: если $f \in \big(\Wpz{p^*}{\ell}(\Om)\big)^*$, то по теореме~\ref{thm:w-dual},
применённой к $\Wpz{p^*}{\ell}(\Om)$, функционал представляется набором
$\underline v \in \big[\Lp{p}(\Om)\big]^N$, а его сужение на $C_0^\infty(\Om)$ — это
в точности распределение $T$ из определения.
\end{proof}

\begin{idea}[title={Почему индекс отрицательный}]
$\Wp{p}{\ell}$ — это, грубо говоря, «первообразные порядка $\ell$»: функции, которые
становятся лучше после интегрирования. $\Wp{p}{-\ell}$ — наоборот, «производные до
порядка $\ell$» от функций из $\Lp{p}$. Элементы $\Wp{p}{-\ell}$ — уже не функции, а
сингулярные распределения.
\end{idea}

\begin{proposition}
$\Lp{p}(\Om)$ всюду плотно в $\Wp{p}{-\ell}(\Om)$.
\end{proposition}

\begin{proof}
\plusplus
\end{proof}
